%% kapa: exam=1 round=31 year=2020 period=1 form=A subject=economics subject_ko=경제학원론 seq=15 number=55
%% source: Q-Net 감정평가사 기출문제 (cst00309) artlSeq=5207219 file=extracted/1008/5207219/A형/A형/2020년도 제31회 감정평가사 1차 시험_1교시_A형.pdf
%% conversion: kordoc 4.24.51 -> tools/gen1.py (auto); review=proofread-sonnet; flags=none
\begin{question}{55}
맥주 시장의 수요함수가 $Q_{D}=100-4P-P_{C}+0.2I$ 일 때, 옳은 것을 모두 고른 것은? (단, $Q_{D}$ 는 맥주 수요량, $P$ 는 맥주 가격, $P_{C}$ 는 치킨 가격, $I$ 는 소득)
\begin{framedtext}
ㄱ. 맥주는 열등재이다. \par
ㄴ. 맥주는 치킨의 보완재이다. \par
ㄷ. 치킨 가격이 인상되면 맥주 수요는 감소한다. \par
\end{framedtext}
\begin{choices}
  \item ㄱ
  \item ㄷ
  \item ㄱ, ㄴ
  \item ㄴ, ㄷ
  \item ㄱ, ㄴ, ㄷ
\end{choices}
\end{question}
\answer{4}
